On s'intéresse à l'évolution de la valeur d'une voiture.

En janvier 2025, la valeur de la voiture est égale à 10\,000 euros.

On suppose qu'ensuite, chaque année, la valeur diminue de 10\,\%.

Pour tout entier naturel $n$, on note $u_n$ la valeur de la voiture, en euros, au mois de janvier de l'année $2025 + n$.

On a donc $u_0 = \np{10000}$.

\begin{enumerate}
	\item \begin{enumerate}
		\item Calculer la valeur de la voiture en janvier 2026.
		\item Justifier que pour tout entier naturel $n$, on a
		\[u_{n+1} = 0,9\, u_n.\]
	\end{enumerate}

	\item En déduire la nature de la suite $(u_n)$ et préciser sa raison.

	\item \begin{minipage}[t]{0.75\linewidth}
		\begin{enumerate}
			\item Pour tout entier naturel $n$, exprimer le terme $u_n$ en fonction de $n$.

			\item En déduire la valeur de la voiture en janvier 2030.

			\emph{On pourra s'aider de l'aide au calcul ci-contre.}
		\end{enumerate}
	\end{minipage}
	\hfill
	\begin{tblr}[t]{|X[3.5cm]|}\hline
		Aide au calcul\\
		$0,9^4 = 0,6561$\\
		$0,9^5 = 0,59049$\\
		$0,9^6 = 0,531441$\\ \hline
	\end{tblr}

	\item \begin{minipage}[t]{0.62\linewidth}
		On considère le programme Python ci-contre.

		Recopier et compléter ce programme afin qu'il renvoie, après exécution, la plus petite valeur de $n$ telle que \quad$u_n < A$,\quad où $A$ est un nombre réel strictement positif.
	\end{minipage}\hfill
	\begin{minipage}[t]{5.5cm}
		\begin{lstlisting}[style=python,numbers=none,xleftmargin=0mm,xrightmargin=0mm]
def seuil(A) :
	N = 0
	U = 10000
	while ......... :
		U = .......
		N = N+1
	return N
\end{lstlisting}
	\end{minipage}

	\item L'exécution de ce programme, pour plusieurs valeurs de $A$, donne les résultats suivants :

\begin{lstlisting}[style=python,xleftmargin=5cm,xrightmargin=5cm,numbers=none]
>>> seuil(7500)
  3
>>> seuil(5000)
  7
>>> seuil(2500)
  14
>>> seuil(2000)
  16
\end{lstlisting}

	\begin{enumerate}
		\item Interpréter, dans le contexte de l'exercice, le résultat obtenu lors de l'appel \texttt{seuil(5000)}.

		\item Déterminer, en expliquant la démarche, l'année à partir de laquelle la voiture aura perdu plus des trois quarts de sa valeur initiale.
	\end{enumerate}
\end{enumerate}
